\section{Discussion}\label{sec:discussion}

\looseness=-1 Our diagnostic exclusion identifies exposure bias as a primary factor in the TF-AR gap but cannot fully rule out KV cache drift (\S\ref{sec:limitations}); PCGrad rules out gradient direction conflict but not magnitude imbalance (7.5--25.7$\times$ NTP dominance; \citealp{chen2018gradnorm}). Despite weak downstream correlation ($r^2$\,=\,0.04), GAR retains diagnostic value by quantifying attention reallocation---a category-dependent effect (+23.8\,pp for content constraints under $\beta$\,=\,3, no change for format/positional). More broadly, the collapse should generalize to any auxiliary objective targeting distributional properties rather than prediction itself \citep{he2019exposure}. Our $\beta$-steering is a \emph{diagnostic}, complementary to inference-time remedies \citep{li2024measuring,hong2025rhea}. These findings resonate with \citet{dou2019attention,dou2022parallel}, who independently observed TF-to-AR attention degradation in NMT. SS's benefit may partly reflect a curriculum effect, though consistent gap compression across seeds and architectures supports the exposure bias interpretation.
